\documentclass[10pt,a4paper]{amsart}
\usepackage[margin=19mm]{geometry}
\usepackage{amsmath,amssymb,mathtools}
\usepackage[scr=rsfs,cal=euler]{mathalfa}
\usepackage{xfrac}
\usepackage[unicode,hidelinks,bookmarks=false]{hyperref}
\newcommand{\ie}{i.e.\ }
\newcommand{\cf}{cf.\ }
\newcommand{\resp}{respectively\ }
\newcommand{\Subsection}{\subsection}
\providecommand{\nobreakdash}{\nobreak\hbox{-}}
\setlength{\emergencystretch}{2em}
\multlinegap0pt
\usepackage{bm}
\usepackage{bbm}
\usepackage{dsfont}
\usepackage{xspace}
\usepackage{cancel}
\usepackage{mathtools}
\usepackage{amsthm}
\makeatletter
\@ifundefined{theorem}{\newtheorem{theorem}{Theorem}}{}
\@ifundefined{lemma}{\newtheorem{lemma}[theorem]{Lemma}}{}
\@ifundefined{proposition}{\newtheorem{proposition}[theorem]{Proposition}}{}
\makeatother
\title[Limiting Absorption Principle for the Helmholtz Equation wit]{Limiting Absorption Principle for the Helmholtz Equation with Sign-Changing Coefficients in Multilayer Spheres}
\author{Wenjing Zhang, Yixian Gao}
\date{Source: arXiv:2605.24826v1 (CC BY 4.0)}
\begin{document}
\maketitle

\noindent\emph{Source and licence.} Limiting Absorption Principle for the Helmholtz Equation with Sign-Changing Coefficients in Multilayer Spheres. Version arXiv:2605.24826v1 (CC BY 4.0), distributed under the Creative Commons Attribution 4.0 International licence, as recorded in the arXiv licence metadata for this exact version and shown on the abstract page https://arxiv.org/abs/2605.24826v1. Full rights notice: licenses/arxiv-2605.24826-CC-BY-4.0.txt.\par\smallskip

\noindent\textit{Editorial scope.} Counted unit: the Proposition whose source label is \texttt{prop:uniform-estimate} in the article (source0.tex lines 428--437, source label \texttt{prop:uniform-estimate}), delivered together with the article's own complete original proof of it (source0.tex lines 440--533, one \texttt{proof} environment of 94 lines). No mathematics is written, repaired or re-derived: statement and proof are the article's own text line for line. Required context retained uncounted: lem:far-field (lines 360--366). Every definition, display and label of those lines is retained unchanged. Omitted from this package and not redistributed: 1--359, 367--427, 438--439, 534--1604. Nothing inside a retained excerpt is omitted or altered: every retained body is delivered exactly as the article's lines are written, so no omission receipt of any kind accompanies this package. Citations: the retained excerpts carry no citation command at all, so no bibliography entry of the article is transcribed and no reference is redistributed. Reference adaptation: none was needed and none was made; every reference inside the delivered unit resolves inside it, so no unresolved reference or citation is produced. Printed slips kept literal: no printed word, symbol, hypothesis or number of the counted statement or of its proof was altered. Layout and class: the article's own class is \texttt{amsart} with packages including amsmath, amssymb, amsthm, mathrsfs, graphicx, cite, newtxtext, newtxmath, cases, color, appendix, enumerate, bm, mathtools; the wrapper is the lane's pinned \texttt{amsart} editorial wrapper at 10pt on A4, which restates the theorem-like environments the retained text needs and adds the article's own macro definitions that the retained text needs (none), with extra wrapper packages (bm, bbm, dsfont, xspace, cancel, mathtools). The delivered numbering is the unit's own. Assets: no retained span contains a figure environment, a graphics-inclusion command, a PDF-page inclusion, a table or a TikZ picture; no image file is copied into this package, the package declares no asset dependency and no image file is redistributed with it. Licence: version arXiv:2605.24826v1 of this article is distributed under the Creative Commons Attribution 4.0 International licence, as recorded in the arXiv licence metadata for this exact version and shown on the abstract page https://arxiv.org/abs/2605.24826v1. A full rights notice is delivered at licenses/arxiv-2605.24826-CC-BY-4.0.txt. Characters: every retained excerpt is pure ASCII, so the delivered engine, XeLaTeX with its default Latin Modern text font, reports no missing character.
\begin{lemma}\label{lem:far-field}
Let $f \in L^2(\mathbb{R}^d)$ have compact support. The radiating solution $u$ to the dissipative Helmholtz equation $(\Delta + \kappa_\sigma^2)u = f$ satisfies the following uniform asymptotic estimate as $R \to \infty$:
\begin{equation*}
    u(x) = O\left( |R|^{-\frac{d-1}{2}} e^{-\operatorname{Im}\kappa_\sigma |R|} \right).
\end{equation*}
The implicit constant depends only on $d$, $\kappa_\sigma$, and the support of $f$.
\end{lemma}

\begin{proposition}\label{prop:uniform-estimate}
Let $f \in L^2(\mathbb{R}^d)$ and let $u_\sigma \in H^1(\mathbb{R}^d)$ denote the unique radiating solution to the dissipative Helmholtz equation
\begin{equation}\label{ge}
    \nabla \cdot (\varepsilon_\sigma^{-1} \nabla u_\sigma) + \omega^2 \mu u_\sigma + \mathrm{i}\sigma u_\sigma = f \quad \text{in } \mathbb{R}^d,
\end{equation}
subject to the outgoing condition at infinity. Then, there exists a constant $C > 0$, independent of both $f$ and $\sigma$, such that the following estimate holds:
\begin{equation*}
    \|u_\sigma\|_{H^1(\mathbb{R}^d)} \le \frac{C}{\sigma} \|f\|_{L^2(\mathbb{R}^d)}.
\end{equation*}
\end{proposition}

\begin{proof}
Testing the governing equation  \eqref{ge} against $\overline{u_\sigma}$ and integrating over the ball $B_R \subset \mathbb{R}^d$, we obtain
\begin{equation*}
    \int_{B_R} \Bigl(\nabla \cdot (\varepsilon_\sigma^{-1} \nabla u_\sigma)\overline{u_\sigma} + \omega^2 \mu |u_\sigma|^2 + \mathrm{i}\sigma |u_\sigma|^2\Bigr)\,\mathrm{d}x = \int_{B_R} f\overline{u_\sigma}\,\mathrm{d}x.
\end{equation*}
Integration by parts gives
\begin{align}\label{eq:integrated-form}
    &\int_{\partial B_R} \varepsilon_\sigma^{-1}\nabla u_\sigma\cdot\vec{n} \overline{u_\sigma}\,\mathrm{d}s - \int_{B_R} \varepsilon_\sigma^{-1} |\nabla u_\sigma|^2\,\mathrm{d}x  \nonumber\\
    &\qquad + \int_{B_R} \omega^2 \mu |u_\sigma|^2\,\mathrm{d}x + \int_{B_R} \mathrm{i}\sigma |u_\sigma|^2\,\mathrm{d}x = \int_{B_R} f\overline{u_\sigma}\,\mathrm{d}x.
\end{align}
In view of Lemma~\ref{lem:far-field}, the solution $u_\sigma$ satisfies the following asymptotic decay as $R \to \infty$:
\begin{equation*}
    u_\sigma = O\!\bigl(R^{-(d-1)/2}e^{-\operatorname{Im}\kappa_\sigma R}\bigr) \quad \text{and} \quad \nabla u_\sigma \cdot \vec{n} = O\!\bigl(R^{-(d-1)/2}e^{-\operatorname{Im}\kappa_\sigma R}\bigr).
\end{equation*}
It follows that the boundary integral vanishes in the limit:
\begin{equation*}
    \left| \int_{\partial B_R} \varepsilon_\sigma^{-1}\nabla u_\sigma\cdot\vec{n} \overline{u_\sigma}\,\mathrm{d}s \right| \leq C R^{d-1} \left( R^{-(d-1)} e^{-2\operatorname{Im}\kappa_\sigma R} \right) \longrightarrow 0 \quad \text{as } R \to \infty.
\end{equation*}
Taking the limit $R \to \infty$ in \eqref{eq:integrated-form}, we obtain
\begin{equation}\label{eq:limit-equation}
    -\int_{\mathbb{R}^d} \varepsilon_\sigma^{-1} |\nabla u_\sigma|^2\,\mathrm{d}x + \int_{\mathbb{R}^d} \omega^2 \mu |u_\sigma|^2\,\mathrm{d}x + \int_{\mathbb{R}^d} \mathrm{i}\sigma |u_\sigma|^2\,\mathrm{d}x = \int_{\mathbb{R}^d} f\overline{u_\sigma}\,\mathrm{d}x.
\end{equation}

Considering the imaginary part of \eqref{eq:limit-equation} and noting the structure of $\varepsilon_\sigma^{-1}$, we find
\begin{equation*}
    \sum_{l} \int_{D_l} \frac{\sigma/\omega}{\varepsilon_l^2 + (\sigma/\omega)^2} |\nabla u_\sigma|^2\,\mathrm{d}x + \int_{\mathbb{R}^d} \sigma |u_\sigma|^2\,\mathrm{d}x = \operatorname{Im} \int_{\mathbb{R}^d} f\overline{u_\sigma}\,\mathrm{d}x.
\end{equation*}
Hence, for each $l$, we have
\begin{align}
\int_{D_l}\dfrac{1}{(\varepsilon_{l}^2+(\sigma/\omega)^2)\omega}|\nabla u_\sigma|^2\mathrm{d}x &\leq \dfrac{1}{\sigma}\left|\int_{\mathbb{R}^d}f\overline{u_\sigma}\mathrm{d}x\right|, \label{eq:l2-grad-l} \\
\|u_\sigma\|_{L^2(\mathbb{R}^d)}^2 &\leq \dfrac{1}{\sigma}\left|\int_{\mathbb{R}^d}f\overline{u_\sigma}\mathrm{d}x\right|. \label{eq:l2-norm-u}
\end{align}
Applying Young's inequality, we have
\begin{equation}\label{eq:gradient-estimate-l}
\|\nabla u_\sigma\|_{L^2(D_l)}^2 \leq \frac{(\varepsilon_l^2 + (\sigma/\omega)^2)\omega}{\sigma}\biggl(\frac{1}{2\eta_l}\|f\|_{L^2(\mathbb{R}^d)}^2 + \frac{\eta_l}{2}\|u_\sigma\|_{L^2(\mathbb{R}^d)}^2\biggr), \quad \forall \eta_l > 0,
\end{equation}
and
\begin{equation}\label{eq:l2-estimate}
\|u_\sigma\|_{L^2(\mathbb{R}^d)}^2 \leq \frac{1}{\sigma}\biggl(\frac{1}{2\eta}\|f\|_{L^2(\mathbb{R}^d)}^2 + \frac{\eta}{2}\|u_\sigma\|_{L^2(\mathbb{R}^d)}^2\biggr), \quad \forall \eta > 0.
\end{equation}
Combining \eqref{eq:gradient-estimate-l} and \eqref{eq:l2-estimate}, we have
\begin{align*}
&\sum_{l} \|\nabla u_\sigma\|_{L^2(D_l)}^2 + \|u_\sigma\|_{L^2(\mathbb{R}^d)}^2 \\
&\qquad \leq \sum_{l} \frac{\omega(\varepsilon_l^2 + (\sigma/\omega)^2)}{\sigma}\biggl(\frac{1}{2\eta_l}\|f\|_{L^2(\mathbb{R}^d)}^2 + \frac{\eta_l}{2}\|u_\sigma\|_{L^2(\mathbb{R}^d)}^2\biggr) + \frac{1}{\sigma}\biggl(\frac{1}{2\eta}\|f\|_{L^2(\mathbb{R}^d)}^2 + \frac{\eta}{2}\|u_\sigma\|_{L^2(\mathbb{R}^d)}^2\biggr).
\end{align*}
Choosing
\[
\eta_l = \frac{\sigma}{\omega(\varepsilon_l^2 + (\sigma/\omega)^2)} \cdot \frac{1}{n-1}, \quad \eta = \frac{\sigma}{2},
\]
 we arrive at the combined estimate
\begin{equation}\label{eq:combined-estimate}
\sum_{l} \|\nabla u_\sigma\|_{L^2(D_l)}^2 + \|u_\sigma\|_{L^2(\mathbb{R}^d)}^2 \leq \frac{C}{\sigma^2}\|f\|_{L^2(\mathbb{R}^d)}^2.
\end{equation}





Next, we take the real part of \eqref{eq:limit-equation} to control the gradient on the remaining domains $D_k$:
\begin{equation*}
    -\sum_{k} \int_{D_k} \varepsilon_k^{-1} |\nabla u_\sigma|^2\,\mathrm{d}x - \sum_{l} \int_{D_l}\dfrac{\varepsilon_l}{\varepsilon_l^2+(\sigma/\omega)^2} |\nabla u_\sigma|^2\,\mathrm{d}x + \int_{\mathbb{R}^d} \omega^2 \mu |u_\sigma|^2\,\mathrm{d}x = \operatorname{Re} \int_{\mathbb{R}^d} f\overline{u_\sigma}\,\mathrm{d}x.
\end{equation*}
Rearranging and invoking \eqref{eq:combined-estimate}, we obtain for each index $k$:
\begin{align*}
\int_{D_k} |\nabla u_\sigma|^2\,\mathrm{d}x &\leq \varepsilon_k\biggl(\sum_{l} \int_{D_l} \frac{|\varepsilon_l|}{\varepsilon_l^2 + (\sigma/\omega)^2} |\nabla u_\sigma|^2\,\mathrm{d}x + \int_{\mathbb{R}^d} \omega^2 \mu |u_\sigma|^2\,\mathrm{d}x + \biggl|\int_{\mathbb{R}^d} f\overline{u_\sigma}\,\mathrm{d}x\biggr|\biggr) \\
&\leq \varepsilon_k\biggl(\sum_{l} \int_{D_l} \frac{|\varepsilon_l|}{\varepsilon_l^2 + (\sigma/\omega)^2} |\nabla u_\sigma|^2\,\mathrm{d}x + \int_{\mathbb{R}^d} \omega^2 \mu^{\mathrm{m}} |u_\sigma|^2\,\mathrm{d}x + \frac{1}{2\eta'}\|f\|_{L^2(\mathbb{R}^d)}^2\\
&\quad+ \frac{\eta'}{2}\|u_\sigma\|_{L^2(\mathbb{R}^d)}^2\biggr),\quad\forall\eta'>0,
\end{align*}
where $\mu^{\mathrm{m}} = \max_i |\mu_i|$. Invoking \eqref{eq:gradient-estimate-l} and \eqref{eq:l2-estimate} and setting the parameters
\[
\eta_l = \frac{\sigma}{\omega\varepsilon_k |\varepsilon_l| (n-1)(n+1)}, \quad \eta = \frac{\sigma}{\omega^2(n+1)\varepsilon_k \mu^{\mathrm{m}}}, \quad \eta' = \frac{1}{2\varepsilon_k(n+1)},
\]
it follows that
\[
\int_{D_k} |\nabla u_\sigma|^2\,\mathrm{d}x \leq \frac{C}{\sigma^2}\|f\|_{L^2(\mathbb{R}^d)}^2 + \frac{1}{n+1}\|u_\sigma\|_{L^2(\mathbb{R}^d)}^2.
\]
Hence
\begin{equation}\label{eq:gradient-estimate-k}
\sum_{k} \|\nabla u_\sigma\|_{L^2(D_k)}^2 \leq \frac{C}{\sigma^2}\|f\|_{L^2(\mathbb{R}^d)}^2 + \frac{1}{2}\|u_\sigma\|_{L^2(\mathbb{R}^d)}^2.
\end{equation}

Combining \eqref{eq:combined-estimate} and \eqref{eq:gradient-estimate-k}, we conclude that
\[
\sum_{k} \|\nabla u_\sigma\|_{L^2(D_k)}^2 + \sum_{l} \|\nabla u_\sigma\|_{L^2(D_l)}^2 + \|u_\sigma\|_{L^2(\mathbb{R}^d)}^2 \leq \frac{C}{\sigma^2}\|f\|_{L^2(\mathbb{R}^d)}^2.
\]
Summing the contributions over the partition $\mathbb{R}^d = \bigcup_i D_i$ and taking the square root, we conclude that
\begin{equation*}
	\|u_\sigma\|_{H^1(\mathbb{R}^d)} \le \frac{C}{\sigma}\, \|f\|_{L^2(\mathbb{R}^d)},
\end{equation*}
which completes the proof.



\end{proof}

\end{document}
